\documentclass[11pt]{article}
\usepackage{amsmath, amssymb}
\usepackage{geometry}
\usepackage{array}
\usepackage{booktabs}
\geometry{margin=0.7in}
\setlength{\parindent}{0pt}
\setlength{\parskip}{0.35em}
\renewcommand{\arraystretch}{1.15}

\newcommand{\cheatsheettitle}[2]{%
\begin{center}
{\LARGE #1}\\[0.4em]
{\large #2}
\end{center}
}


\begin{document}
\cheatsheettitle{Algebra I Statistics and Probability Cheatsheet}{Module 11}

\section*{Measures of Center}
\[
\text{mean}=\frac{\text{sum of data values}}{\text{number of data values}}
\]
Median is the middle value after ordering the data. Mode is the most frequent value.

\section*{Measures of Spread}
\[
\text{range}=\max-\min
\]
\[
\text{IQR}=Q_3-Q_1
\]
Five-number summary:
\[
\min,\quad Q_1,\quad \text{median},\quad Q_3,\quad \max
\]

\section*{Box Plots}
\begin{itemize}
  \item Box runs from $Q_1$ to $Q_3$.
  \item Median line is inside the box.
  \item Whiskers extend to minimum and maximum unless outliers are separated.
\end{itemize}
Outlier rule:
\[
x<Q_1-1.5(\mathrm{IQR})\quad \text{or}\quad x>Q_3+1.5(\mathrm{IQR})
\]

\section*{Two-Way Tables}
Rows and columns show categories. Joint frequencies are inside the table. Marginal frequencies are row or column totals.
\[
P(A)=\frac{\text{count in }A}{\text{total count}}
\]
\[
P(A\mid B)=\frac{\text{count in both }A\text{ and }B}{\text{count in }B}
\]

\section*{Basic Probability}
\[
P(\text{event})=\frac{\text{favorable outcomes}}{\text{total outcomes}}
\]
\[
0\le P(A)\le1
\]
Complement:
\[
P(\text{not }A)=1-P(A)
\]

\section*{Compound Probability}
Addition rule:
\[
P(A\text{ or }B)=P(A)+P(B)-P(A\text{ and }B)
\]
Independent multiplication rule:
\[
P(A\text{ and }B)=P(A)P(B)
\]

\section*{Odds}
\[
\text{odds in favor}=\frac{\text{favorable outcomes}}{\text{unfavorable outcomes}}
\]
\[
\text{odds against}=\frac{\text{unfavorable outcomes}}{\text{favorable outcomes}}
\]

\section*{Quick Checks}
\begin{itemize}
  \item Order data before finding median or quartiles.
  \item Use totals carefully in two-way tables.
  \item Probabilities can be fractions, decimals, or percents.
  \item ``Or'' includes overlap unless the events are mutually exclusive.
\end{itemize}

\end{document}
